\section{Defense D10: Contextual Integrity Verification}
\label{sec:civ}

This appendix describes \civ{} (D10), our novel multi-signal defense based on Nissenbaum's theory of contextual integrity \citep{nissenbaum2004privacy}.
The key insight is that indirect prompt injection violates the \emph{contextual integrity} of the agent's information flow: tool calls triggered by injected instructions reference entities, use tools, or exhibit patterns that are inconsistent with the user's original goal.

\subsection{Motivation: Contextual Integrity for Agent Security}

Nissenbaum's theory defines appropriate information flow through \emph{informational norms}: rules governing (1)~the actors involved, (2)~the type of information, and (3)~the transmission principle.
We adapt this framework to the agent setting:
\begin{itemize}
    \item \textbf{Actors}: the user (goal provider), the untrusted data source (content provider), and the agent (executor).
    \item \textbf{Information types}: entities referenced in tool-call parameters (email addresses, file paths, URLs).
    \item \textbf{Transmission principle}: tool calls should serve the user's goal, not instructions embedded in untrusted content.
\end{itemize}

A violation of contextual integrity occurs when the agent makes a tool call whose parameters or purpose trace back to the untrusted content rather than the user's goal.

\subsection{CIV Design: From v1 to v2}
\label{sec:civ_design}

We developed \civ{} iteratively.
\textbf{CIV~v1} implemented three signals: entity provenance (Signal~A), a keyword-based tool-affinity fingerprint (Signal~B), and a heuristic counterfactual reasoning check (Signal~C).
Ablation of v1 (\S\ref{sec:ablation}) revealed two critical weaknesses: (1)~the counterfactual heuristic is too permissive---removing it \emph{improves} ASR by 53\%; (2)~keyword-based tool affinity lacks semantic generalization.

\textbf{CIV~v2} (reported in the main table) incorporates these findings:
\begin{itemize}
    \item Signal~B is upgraded from keyword lookup to \textbf{embedding-based semantic compatibility}, computing cosine similarity between the goal and tool-description embeddings via a pre-trained language model, with keyword matching as a fallback.
    \item Signal~C replaces counterfactual reasoning with \textbf{plan deviation detection}: an expected tool set is inferred from goal keywords and tool co-occurrence statistics, and tools outside this plan are penalized.
    \item A \textbf{risk-stratified gating} layer distinguishes read-only operations (which cannot cause data exfiltration) from side-effect operations, applying a lenient threshold to the former to reduce false positives on multi-step benign tasks.
\end{itemize}

\subsection{Three-Signal Verification (CIV v2)}
\label{sec:three_signals}

\civ{} v2 evaluates each proposed tool call using three complementary signals fused into a combined integrity score.

\paragraph{Signal A: Entity Provenance with Session Accumulation.}
The provenance signal checks whether tool-call parameters reference entities originating from untrusted content rather than the user's goal.

Given goal $g$ and untrusted content $c$, we extract entity sets $E_g$ and $E_c$ using pattern matching for email addresses, file paths, URLs, quoted strings, and identifiers.
Let $E_p$ be the entities in the proposed tool call's parameters, and $T$ be the \emph{session trust set}---initially $E_g$, extended with entities seen in all previously allowed tool calls.
\begin{equation}
\small
    s_{\text{prov}} {=} \begin{cases}
        1.0 & E_p {=} \emptyset \text{ or } E_p \cap (E_c {\setminus} T) {=} \emptyset \\[4pt]
        \frac{|E_p \cap T|}{|E_p \cap T| {+} |E_p \cap (E_c {\setminus} T)|} & \text{otherwise}
    \end{cases}
\end{equation}
Session accumulation prevents cascading false positives in multi-tool tasks: once an entity is seen in an approved tool call, it is admitted into the trust set for subsequent calls.

\paragraph{Signal B: Embedding-Based Semantic Compatibility.}
The compatibility signal evaluates whether the tool call is semantically consistent with the user's goal.
Each tool $t$ is associated with a natural-language description $d_t$ (e.g., \texttt{send\_email}: ``Send, forward, or reply to email messages'').
The goal-tool compatibility score is:
\begin{equation}
    s_{\text{compat}} = \cos\bigl(\text{embed}(g),\ \text{embed}(d_t)\bigr)
\end{equation}
normalized to $[0, 1]$.
When tool-call history $h$ is available, the score is blended with a co-occurrence affinity term:
\begin{equation}
\small
    s_{\text{compat}} = 0.6 \cdot \cos\!\bigl(\text{embed}(g), \text{embed}(d_t)\bigr) + 0.4 \cdot \max_{t_i \in h} A(t_i, t)
\end{equation}
where $A(t_i, t)$ is the empirical co-occurrence affinity between tools $t_i$ and $t$.
A keyword-based fallback is used when the embedding API is unavailable.

\paragraph{Signal C: Plan Deviation Detection.}
The plan deviation signal checks whether the proposed tool belongs to the expected tool plan inferred from the user's goal.
We extract an \emph{expected plan} $P_g$ by mapping goal keywords to likely tools using a curated keyword-to-tool implication table.
\begin{equation}
    s_{\text{plan}} = \begin{cases}
        1.0 & \text{if } t \in P_g \\
        0.6 & \text{if } \exists\, t_i \in P_g \cup h \text{ s.t. } A(t_i, t) \geq 0.2 \\
        0.2 & \text{otherwise}
    \end{cases}
\end{equation}
Unlike the counterfactual heuristic in v1---which estimated whether the tool call would still occur absent untrusted content and was found to be net-harmful---plan deviation detection uses only the user goal and co-occurrence priors, avoiding reliance on the untrusted content itself.

\subsection{Signal Fusion and Risk-Stratified Gating}

The three signals are fused with weights $(w_{\text{prov}}, w_{\text{compat}}, w_{\text{plan}}) = (0.45, 0.30, 0.25)$:
\begin{equation}
    s_{\text{combined}} = w_{\text{prov}} \cdot s_{\text{prov}} + w_{\text{compat}} \cdot s_{\text{compat}} + w_{\text{plan}} \cdot s_{\text{plan}}
\end{equation}

\textbf{Risk-stratified gating} applies different thresholds based on whether the tool has side effects:
\begin{itemize}
    \item \textbf{Read-only tools} (e.g., \texttt{search\_web}, \texttt{read\_email}): lenient path with threshold $\tau_r = 0.30$ and equal weighting of plan and compatibility signals only. Read-only operations cannot exfiltrate data by themselves and are relaxed to reduce false positives in multi-step benign tasks.
    \item \textbf{Side-effect tools} (e.g., \texttt{send\_email}, \texttt{write\_file}): strict path with threshold $\tau_w = 0.45$, increased by $+0.20$ for critical-risk tools and $+0.10$ for high-risk tools.
\end{itemize}

\subsection{Ablation Findings and Design Rationale}
\label{sec:civ_ablation_findings}

Our ablation study of CIV~v1 (Table~\ref{tab:civ_ablation}) directly motivated the v2 design:
\begin{enumerate}
    \item \textbf{Provenance is the core signal}: Removing provenance causes the largest ASR increase (+115\%), confirming entity provenance as the dominant attack detector.
    Single-signal provenance achieves 9.7\% ASR---nearly matching the full v1 system---and forms the backbone of CIV~v2.
    \item \textbf{Counterfactual reasoning is counterproductive}: Removing counterfactual from v1 \emph{improves} security ($-53\%$ ASR), motivating its replacement with plan deviation detection in v2.
    \item \textbf{Tool affinity supports utility}: The fingerprint signal contributes most to maintaining BSR ($-12\%$ when removed), which is preserved in v2 via embedding-based compatibility.
\end{enumerate}

CIV~v2 achieves 8.9\% ASR and 82.1\% BSR on the 250-case core benchmark, compared to 12.3\% ASR and 85.3\% BSR for CIV~v1-full and 5.8\% ASR / 82.1\% BSR for v1 without counterfactual.
The v2 design thus balances the security improvement from removing counterfactual against the utility cost of replacing the permissive heuristic with structured plan deviation.

\subsection{Limitations}

CIV's entity provenance signal struggles with multi-step benign workflows that legitimately move information across tools (see \S\ref{sec:limitations}).
The plan deviation model uses a curated keyword-to-tool implication table and empirical co-occurrence statistics; learning these priors from agent interaction logs would improve generalization to new tool ecosystems.
Future work should investigate learned tool affinity models and principled counterfactual evaluation using LLM-based plan verification.

